\section{Entity Extraction}
\label{sec:information_extraction}

Capturing the meaning of text is a central goal of \acl{nlp}~\cite{entityorientedsearch2018balog}, and the automatic extraction of semantic information from text---such as mentions of named entities, relationships, and attributes---is known as \emph{\acl{ie}}. This has been the objective of major competitions, including the Message Understanding Conference (MUC)~\cite{muc_lessons_1998,muc6_overview_1995,muc7_overview_1998}, the Computational Natural Language Learning (CoNLL) 2003 shared task~\cite{conll2003}, and the Automatic Content Extraction (ACE) program~\cite{ace_2004}.

% Combining the higher-quality information from \ac{kb} with the broader coverage of unstructured data has motivated the research in \emph{\acl{ie}} since three decades ago~\cite{infoextraction2008,entityorientedsearch2018balog}.
\begin{definition}[Information Extraction]
``\Acf{ie} refers to the automatic extraction of structured information such as entities, relationships between entities, and attributes describing entities from unstructured sources''~\cite{infoextraction2008}.
\end{definition}
The use of \ac{ie} techniques is generally aimed at either \emph{\acl{sta}} or \emph{\acl{kbp}}~\cite{entityorientedsearch2018balog}.
In the former case, annotating text with information from a \ac{kb} enriches the text with additional semantics, enabling advanced information access functionalities, such as entity-based search and filtering via faceting~\cite{entityorientedsearch2018balog,ARMENTANO2014_nlp_faceted,hearst_2006_faceted}.
In the latter case, \ac{ie} is used to populate or update a \ac{kb} with information extracted from unstructured sources to enhance the \ac{kb}'s coverage and keep it up-to-date.
In this work, we focus \acf{ee} concentrating on entities, without considering broader \ac{ie} tasks such as relation extraction. 

In order, we define and describe the literature of \acl{er} (\Cref{sec:rw:ner}), \acl{nel} (\Cref{sec:rw:nel}), NIL prediction and clustering (\Cref{sec:incrementalel}), \ac{ee} in the legal domain (\Cref{sec:sota:ee:legal}), and finally we describe \ac{ee} resources for the Italian language (\Cref{sec:sota:ee:ita}).

% \todo{use tac kbp for introducing ner nel and co}



% \textbf{extracting knowledge} and its difference with bare information extraction. Integrating knowledge (e.g. NEL), semantic web,... inferring relations,...

% Describe and give the definitions always continuing the examples.

% Knowledge in general? types of knowledge? Knowledge Graphs, Knowledge representation
% Entity Recognition, Linking and other tasks

% How to extract Knowledge

% How can machine access Knowledge?

% different ways for accessing external knowledge

% maybe move SOTA in the respective chapters
